% Drawing 76/01-5 -- folder 76, batch 01, archive page 9 (PDF page 10).
% Transcription: au bas de la page, à gauche, sept petites figures sans lettres, exemples de 1-cartes : une boucle fermée sur un point ; un segment entre deux gros sommets ; un segment entre un gros sommet et un petit ; un « Y » dont deux bouts sont des sommets cerclés et le troisième un sommet plein ; un court segment isolé terminé par un point ; un disque dont le bord est une boucle issue d'un point, traversé par un arc issu du même point ; un sommet d'où partent quatre arêtes, deux épaisses entre lesquelles sont tracées des hachures.
% Kind: seven small figures: examples of 1-maps (loops, segments, a Y, a vertex with edges). Method 2 (schematic redraw from the facsimile). Status: ok.
% Notes: As on the page, each figure through its own computed points: (1) a loop closed on a full point; (2) a segment between two heavy points, with a small tick at its middle; (3) a segment between a heavy point and a smaller one; (4) a Y: three arms from one centre, two ending in ringed points and the lower one in a full point; (5) a short isolated segment with a point at its upper end; (6) a full point with four strokes: an upper arc from the point to a tip at the top, an inner arc from just right of the point up to the same tip, and, from a small hook beside that tip, the outer loop round the right and the bottom back to the point and a second inner arc down to just right of the point (a disc bordered by a loop, cut by two inner arcs; the top is not closed, the tip and the hook stand 1.5 mm apart); (7) a full point with four edges: two heavy ones to the right, a longer thin one to the left ending in a small hook, a short thin one to the lower left, and six short hatching strokes between the two heavy edges. Left out: the formulas written under and beside them, which the transcription gives. The transcription's note says six figures and lists seven; it gives three edges to the last vertex where the page has four (the fourth is short and thin). Overlaid on the 400 dpi colour crop on 2026-11-12: figure 6 had one closed curve and one inner arc, the page has the four strokes above (the second inner arc was missing); figure 7 has six hatching strokes, not four; both corrected. Redrawn 2026-11-12 (issue 128): to check again.
% Coordinates: 1 unit = 1 mm = 1/100 of the crop width (bbox in page fractions: [0.1, 0.61, 0.5, 0.85]).
% Redrawn by Claude (claude-fable-5-2), issue #128, 2026-11-12; not checked by a human; not compiled (no TeX in the session).
% Revised by Claude (claude-opus-5-5) 2026-11-12 (issue #128, overlay check on the 400 dpi colour facsimile): figure 6 redrawn as its four strokes (the second inner arc was missing and the top is not closed: tip and hook stand apart); figure 7 has six hatching strokes, not four. Not compiled. Status: ok.
% Rights: reproduces part of the work -- behind RENDER_TRANSCRIPTS, never published by itself.
\documentclass[tikz,border=3pt]{standalone}
\usepackage{amsmath,amssymb}
\usetikzlibrary{patterns,arrows.meta}
\tikzset{pencil/.style={black!45}, struck/.style={black!35}, mpt/.style={circle,draw,inner sep=1.3pt}, pt/.style={circle,fill,inner sep=0.9pt}, opt/.style={circle,draw,fill=white,inner sep=1.1pt}}
\begin{document}
\begin{tikzpicture}[x=1mm,y=1mm,line cap=round,line join=round,font=\scriptsize]
\draw plot[smooth cycle, tension=0.55] coordinates {(10.8,-17.8) (14,-13) (19,-9.6) (23,-9.8) (24.5,-12.5) (22,-15.5) (17,-18) (13,-18.6)};
\node[pt] at (10.8,-17.8) {};
\draw (26.5,-23.5) -- (40,-16.5);
\node[pt] at (26.5,-23.5) {};
\node[pt] at (40,-16.5) {};
\draw (32.9,-19.6) -- (33.5,-20.8);
\draw (48.5,-13.5) -- (71,-12);
\node[pt] at (48.5,-13.5) {};
\node[pt] at (71,-12) {};
\draw (53.3,-33.7) -- (46,-26.8);
\draw (53.3,-33.7) -- (63,-28);
\draw (53.3,-33.7) -- (51.8,-42);
\node[mpt] at (46,-26.8) {};
\node[mpt] at (63,-28) {};
\node[pt] at (51.8,-42) {};
\draw (80.5,-42) -- (85.3,-36.6);
\node[pt] at (85.3,-36.6) {};
\draw plot[smooth, tension=0.55] coordinates {(38.4,-58.3) (41,-54.7) (45.6,-51.5) (50.9,-50.2)};
\draw plot[smooth, tension=0.55] coordinates {(40.7,-58.3) (45.6,-56.5) (49.1,-54.2) (51.2,-51.5)};
\draw plot[smooth, tension=0.55] coordinates {(52.3,-50.6) (52.6,-50) (53.2,-50.2) (55.3,-52.4) (55.8,-56.5) (54.7,-60.3) (51.5,-62.6) (46.8,-63) (42.2,-61.7) (39.2,-59.7) (38.4,-58.3)};
\draw plot[smooth, tension=0.55] coordinates {(53.2,-51) (52.9,-54.7) (50.9,-57.9) (45.6,-59.3) (39.8,-59.3)};
\node[pt] at (38.4,-58.3) {};
\draw[very thick] (80.7,-62.7) -- (87.8,-54.5);
\draw[very thick] (80.7,-62.7) -- (89,-70.5);
\draw plot[smooth, tension=0.55] coordinates {(80.7,-62.7) (75,-61.8) (70.5,-61.5) (69.3,-60.6)};
\draw (80.7,-62.7) -- (77.8,-68.5);
\node[pt] at (80.7,-62.7) {};
\draw (85.5,-58.1) -- (88.5,-58.2);
\draw (85.3,-59.4) -- (87.3,-59.5);
\draw (84.4,-61.3) -- (86.1,-61.3);
\draw (84.1,-63.1) -- (85.8,-62.9);
\draw (84.7,-64.9) -- (86.7,-64.9);
\draw (85,-66.8) -- (87.9,-66.6);
\end{tikzpicture}
\end{document}
